\phantomsection\section*{ЗАКЛЮЧЕНИЕ}\addcontentsline{toc}{section}{ЗАКЛЮЧЕНИЕ}

В ходе выполнения научно-исследовательской работы было проведено комплексное исследование автоматизированных систем
для проверки оформления чертежей.
В соответствии с поставленной целью и задачами были получены следующие результаты:

\begin{enumerate}
\item Введены определения базовых понятий в исследуемой области.
\item Проведён анализ предметной области, классифицированы виды ошибок, приведен стандартный вид системы автоматизированный проверки чертежей.
\item Проведён анализ существующих программных реализаций. Исследованы 30 систем для нахождения ошибок оформления в чертежах. Выявлено, что современные решения все больше и больше используют нейронные сети.
\item Сформулированы критерии сравнения средств автоматической проверки чертежей на ошибки оформления для оценки эффективности систем.
\item Выполнен сравнительный анализ методов, построена сравнительная таблица существующих программных средств для автоматизированной проверки чертежей на ошибки оформления. Установлено, что ни один из методов не исключает наличия ошибок.
    \begin{itemize}
        \item методы, основанные на алгоритмической проверке формальных правил, демонстрируют высокую точность, воспроизводимость и эффективность при выявлении типовых ошибок, явно определённых в нормативной документации. Такие подходы обеспечивают детерминированный результат, легко поддаются верификации и интеграции в существующие САПР;
        \item методы, основанные на искусственном интеллекте обладают способностью к семантическому анализу: они могут интерпретировать смысл изображённых объектов, распознавать нестандартные или сложные геометрические конструкции, а также выявлять логические несоответствия, не описанные явно в формальных правилах. Однако такие системы уступают алгоритмическим методам в надёжности и точности при проверке строгих нормативных требований и страдают от недостатка качественно размеченных данных.
    \end{itemize}
\end{enumerate}

Таким образом, цель работы достигнута в полном объёме. Проведенное исследование показало, что наиболее перспективный способ для создания эффективной
системы автоматизированный проверки чертежей на ошибки оформления --- комплексный подход, заключающийся в использовании обоих методов.

Полученные результаты могут быть использованы для интеграции в системы автоматического проектирования, что позволит сэкономить время прохождения нормоконтроля.
\clearpage
